\section{噪声预测与 Score Matching 的联系}

\subsection{从噪声预测到 Score Function}

讲者在课程结尾预告了一个重要的联系：噪声预测网络 $\epsilon_\theta(x_t, t)$ 与 score-based 模型中的 score function 有着深刻的对应关系。

Score function 定义为对数概率密度对数据的梯度：

$$s_\theta(x_t, t) = \nabla_{x_t} \log p_t(x_t)$$

可以证明，噪声预测网络与 score function 之间存在如下关系：

$$\epsilon_\theta(x_t, t) \approx -\sqrt{1 - \bar{\alpha}_t}\, \nabla_{x_t} \log q_t(x_t)$$

\begin{knowledgebox}{DDPM 与 Score-Based 模型的统一}
这一联系揭示了两个看似不同的生成模型框架——denoising diffusion probabilistic models（DDPM）和 score-based generative models（SGM）——本质上是等价的：
\begin{itemize}
    \item DDPM 学习预测噪声 $\epsilon$
    \item SGM 学习预测 score function $\nabla_x \log p(x)$
    \item 两者之间仅差一个与时间步相关的缩放因子
\end{itemize}
这种等价性由 Song et al.（2021）在 Score SDE 框架中给出了统一的连续时间描述。
\end{knowledgebox}

\subsection{直觉理解}

噪声预测可以被理解为回答这个问题：``$x_t$ 中有多少噪声？''而 score function 回答的是：``$x_t$ 应该往哪个方向移动才能增加其概率密度？''这两个问题在扩散过程的上下文中本质上是等价的——去除噪声的方向正是增加概率密度的方向。

\subsection{本章小结}

噪声预测与 score function 之间存在精确的数学对应关系，这建立了 DDPM 和 score-based 模型之间的桥梁，为后续的 Score SDE 统一框架奠定了基础。

%% ========== 总结与延伸 ==========

